\clearpage
\SourceLeaf{612}
\ArithmeticHead{583}{R}{INHALT.}
\begingroup
\newcommand{\ContentsRow}[5]{\par\noindent\hangindent=#1\hspace*{#1}#2\hspace{.4em}\mbox{\footnotesize #3}\hfill\mbox{#4\quad #5}\par}
\ContentsRow{2em}{Erläuterungen zum Abschnitt VII}{.\,.\,.\,.\,.\,.\,.\,.\,.\,.\,.\,.\,.\,.\,.\,.\,.\,.\,.}{Seite}{438}
\ContentsRow{2em}{Erläuterungen zum Briefwechsel. Geschichtliches zu den Abschnitten V, VI, VII}{.\,.}{—}{442}
\begin{center}\rule{.13\linewidth}{.4pt}\end{center}
\begin{center}GEOMETRIE.\end{center}
\noindent\emph{Nachlass.}\par
\ContentsRow{1em}{Ansätze zur transzendenten Trigonometrie}{.\,.\,.\,.\,.\,.\,.\,.\,.\,.\,.\,.\,.\,.\,.\,.\,.\,.}{—}{451}
\ContentsRow{2em}{Bemerkungen}{.\,.\,.\,.\,.\,.\,.\,.\,.\,.\,.\,.\,.\,.\,.\,.\,.\,.\,.\,.\,.\,.\,.\,.\,.}{—}{453}
\ContentsRow{1em}{Zur sphärischen Trigonometrie (aus \textsc{Wittsteins} \emph{Lehrbuch der Elementar-Mathematik}, 1862)}{}{—}{457}
\ContentsRow{2em}{Bemerkung}{.\,.\,.\,.\,.\,.\,.\,.\,.\,.\,.\,.\,.\,.\,.\,.\,.\,.\,.\,.\,.\,.\,.\,.\,.\,.}{—}{458}
\par\medskip\noindent\emph{Briefwechsel:}\par
\ContentsRow{1em}{1.\ \textsc{Schumacher} an \textsc{Gauss}, 1836 März 19}{.\,.\,.\,.\,.\,.\,.\,.\,.\,.\,.\,.\,.\,.\,.\,.\,.}{—}{459}
\ContentsRow{1em}{2.\ \textsc{Gauss} an \textsc{Schumacher}, 1836 März 21}{.\,.\,.\,.\,.\,.\,.\,.\,.\,.\,.\,.\,.\,.\,.\,.\,.}{—}{461}
\ContentsRow{1em}{3.\ \textsc{Schumacher} an \textsc{Gauss}, 1836 April 2}{.\,.\,.\,.\,.\,.\,.\,.\,.\,.\,.\,.\,.\,.\,.\,.\,.\,.}{—}{462}
\ContentsRow{1em}{4.\ \textsc{Schumacher} an \textsc{Gauss}, 1836 April 7}{.\,.\,.\,.\,.\,.\,.\,.\,.\,.\,.\,.\,.\,.\,.\,.\,.\,.}{—}{462}
\ContentsRow{1em}{5.\ \textsc{Gauss} an \textsc{Schumacher}, 1836 April 13}{.\,.\,.\,.\,.\,.\,.\,.\,.\,.\,.\,.\,.\,.\,.\,.\,.}{—}{464}
\ContentsRow{1em}{6.\ \textsc{Schumacher} an \textsc{Gauss}, 1836 April 19}{.\,.\,.\,.\,.\,.\,.\,.\,.\,.\,.\,.\,.\,.\,.\,.\,.}{—}{466}
\ContentsRow{1em}{7.\ \textsc{Gauss} an \textsc{Schumacher}, 1836.\ April 21}{.\,.\,.\,.\,.\,.\,.\,.\,.\,.\,.\,.\,.\,.\,.\,.\,.}{—}{467}
\ContentsRow{1em}{Bemerkung}{.\,.\,.\,.\,.\,.\,.\,.\,.\,.\,.\,.\,.\,.\,.\,.\,.\,.\,.\,.\,.\,.\,.\,.\,.\,.\,.\,.}{—}{468}
\par\medskip\noindent\textls[120]{Mannigfaltigkeiten von} $n$ \textls[120]{Dimensionen.}\par
\noindent\hspace*{1em}I.\quad Über die Ermittelung des kleinsten Werthes für das Maaß des Zwanges.\par
\ContentsRow{2em}{(Aus der \emph{Inauguraldissertation} von A. Ritter, Göttingen 1853, S.\ 20—23)}{.\,.\,.\,.\,.}{—}{469}
\noindent\hangindent=2em\hspace*{1em}II.\quad Bestimmung des kleinsten Werthes der Summe $x_1^2+x_2^2+\mathord{\cdot}\,\mathord{\cdot}\,\mathord{\cdot}+x_n^2=R^2$ für $m$ gegebene Ungleichungen $u\geqq0$.\par
\ContentsRow{2em}{(Aus \textsc{Gauss}’ Vorlesung \emph{Über die Methode der kleinsten Quadrate}, Wintersem.\ 1850/51)}{}{—}{473}
\ContentsRow{3em}{Bemerkung}{.\,.\,.\,.\,.\,.\,.\,.\,.\,.\,.\,.\,.\,.\,.\,.\,.\,.\,.\,.\,.\,.\,.\,.\,.}{—}{481}
\begin{center}\rule{.13\linewidth}{.4pt}\end{center}
\begin{center}\begin{minipage}[c]{.78\linewidth}\centering NACHBILDUNG DES TAGEBUCHS (NOTIZENJOURNALS) VON C. F. GAUSS,\par 1796 MART. 30—1814 IUL. 9.\end{minipage}\(\left.\vphantom{\begin{array}{c}1\\2\\3\end{array}}\right\}\)\begin{minipage}[c]{.14\linewidth}zwischen-\\geheftet.\end{minipage}\end{center}
\begin{center}\rule{.13\linewidth}{.4pt}\end{center}
\begin{center}ABDRUCK DES TAGEBUCHS (NOTIZENJOURNALS) MIT ERLAEUTERUNGEN.\end{center}
\ContentsRow{0pt}{\emph{Vorbemerkungen zum ersten Abdruck des Tagebuchs}}{.\,.\,.\,.\,.\,.\,.\,.\,.\,.\,.\,.\,.\,.\,.\,.}{—}{485}
\ContentsRow{0pt}{\emph{Vorbemerkung zu dem hier folgenden Abdruck des Tagebuchs}}{.\,.\,.\,.\,.\,.\,.\,.\,.\,.\,.\,.\,.}{—}{487}
\par\medskip
\ContentsRow{0pt}{1.\quad Siebzehntheilung des Kreises, 1796 März 30}{.\,.\,.\,.\,.\,.\,.\,.\,.\,.\,.\,.\,.\,.\,.\,.\,.\,.}{—}{488}
\ContentsRow{0pt}{2.\quad Zum ersten Beweise des Reziprozitätsgesetzes der quadratischen Reste, 1796 April 8}{.\,.\,.}{—}{489}
\ContentsRow{0pt}{3.\quad Zur Kreisteilung, 1796 April 12}{.\,.\,.\,.\,.\,.\,.\,.\,.\,.\,.\,.\,.\,.\,.\,.\,.\,.\,.\,.\,.\,.}{—}{489}
\ContentsRow{0pt}{4.\quad Allgemeiner quadratischer Restcharakter, 1796 April 29}{.\,.\,.\,.\,.\,.\,.\,.\,.\,.\,.\,.\,.}{—}{489}
\ContentsRow{0pt}{5.\quad \textsc{Goldbach}scher Satz, 1796 Mai 14}{.\,.\,.\,.\,.\,.\,.\,.\,.\,.\,.\,.\,.\,.\,.\,.\,.\,.\,.\,.\,.}{—}{490}
\ContentsRow{0pt}{6.\quad Potenzsummen der Wurzeln einer Gleichung, 1796 Mai 23}{.\,.\,.\,.\,.\,.\,.\,.\,.\,.\,.\,.}{—}{490}
\endgroup
